\section{Sistematika Penulisan}
\label{sec:sistematikapenulisan}

Laporan ini disusun secara sistematis ke dalam lima bab untuk memberikan gambaran yang jelas dan terstruktur mengenai seluruh proses pengerjaan proyek. Berikut adalah sistematika penulisannya:
\begin{itemize}
    \item \textbf{BAB I: PENDAHULUAN}\\
    Bab ini berisi latar belakang yang menjadi dasar pemikiran proyek, batasan masalah yang mendefinisikan ruang lingkup pengerjaan, tujuan yang ingin dicapai, serta sistematika penulisan laporan.
    \item \textbf{BAB II: TINJAUAN UMUM INSTITUSI}\\
    Bab ini menguraikan profil singkat institusi tempat praktik kerja lapangan dilaksanakan, yaitu Bapenda Kota Surabaya, beserta sejarah singkat, visi, dan misi perusahaan.
    \item \textbf{BAB III: TINJAUAN PUSTAKA}\\
    Bab ini membahas landasan teori dan konsep-konsep dasar yang digunakan dalam pengembangan proyek, meliputi penjelasan mengenai \emph{E-Government}, Pendapatan Asli Daerah, ASP.NET Core MVC, Entity Framework Core, Oracle, dan lainnya.
    \item \textbf{BAB IV: PEMBAHASAN}\\
    Bab ini merupakan inti dari laporan yang menjelaskan secara rinci proses pengembangan sistem, mulai dari tinjauan kasus, analisis kebutuhan fungsional, hingga pembahasan solusi dan implementasi teknis dari fitur-fitur yang telah dikembangkan serta hasil dan dampaknya.
    \item \textbf{BAB V: PENUTUP}\\
    Bab ini berisi kesimpulan dari keseluruhan hasil proyek pengembangan MonPD Reborn serta saran-saran untuk pengembangan lebih lanjut di masa mendatang.
\end{itemize}
